\input{../tex-inputs/latex-leseansicht-vorspann}

\section[Gustav Schwarzkopf an Arthur Schnitzler, 6. 3. 1908]{L04270 Gustav Schwarzkopf an Arthur Schnitzler, 6. 3. 1908}
\nopagebreak\mylabel{L04270v}
\rehead{ }\normalsize\beginnumbering\briefempfaengerindex{Schnitzler, Arthur@\textsc{Schnitzler, Arthur}!zzzSchwarzkopf, Gustav@\emph{von Gustav Schwarzkopf}!1908-03-062@{6. 3. 1908}|(be}
\toendnotes[C]{\smallbreak\pagebreak[2]}
\correspDesc{Versand  durch Gustav Schwarzkopf am 6. 3. 1908 in Wien
\newline{}Erhalt  durch Arthur Schnitzler im Zeitraum [6. 3. 1908 – 9. 3. 1908?] in Wien}\toendnotes[C]{\smallbreak}
\Standort{CUL, Schnitzler, B 96a.}
\physDesc{Postkarte, 383 Zeichen
\newline{}Handschrift: schwarze Tinte, deutsche Kurrent
\newline{}Versand: Stempel: »\nobreak{}\oindex{I., Innere Stadt@\textbf{I., Innere Stadt}, \emph{Verwaltungsgebiet}|pwk}1/1 Wien 9, 6 III 08, 8\nobreak{}«.  
\newline{}Schnitzler: mit Bleistift Vermerk: »\textsc{Schwarzkopf}« }\toendnotes[C]{\smallbreak}\pstart{}{\pb}I. Tiefer Graben 17\oindex{Wien@\textbf{Wien}!I., Innere Stadt@\textbf{I., Innere Stadt}!Tiefer Graben 17@\textbf{Tiefer Graben 17}, \emph{Wohngebäude}|pw}.\pend{}{\bigskip}\pstart{}Herrn\pend{}\pstart{}\textsc{D\textsuperscript{r} Arthur Schnitzler}\pend{}\pstart{}\textsc{Wien\oindex{Wien@\textbf{Wien}, \emph{Verwaltungsgebiet}|pw}}\pend{}\pstart{}XVIII Spöttelgaſſe 7\oindex{Wien@\textbf{Wien}!XVIII., Währing@\textbf{XVIII., Währing}!Edmund-Weiß-Gasse 7@\textbf{Edmund-Weiß-Gasse 7}, \emph{Wohngebäude}|pw}.\pend{}{\bigskip}\vspace{1em}
\pstart
           \raggedleft{}{\pb}6. III 08\pend
           
\pstart{}Lieber Arthur!\pend\vspace{0.5em}
\pstart
           das{ }ſchöne Wetter hatte mich veranlaßt  früher \introOben{}fortzugehen\introOben{} und als Ihre
               \label{K_L04270-1v}\edtext{pneumatiſche Karte}{\lemma{\textnormal{\emph{pneumatische Karte}}}\Cendnote{\textnormal{XXXX Auszeichnungsfehler: Dokument L04031 nicht gefunden. Pneumatisch: mit Rohrpost.}}}\label{K_L04270-1} kam, war ich nicht mehr, zu Hauſe. Als ist dann gegen 9 Uhr Ihre
               Karte vorfand, war es natürlich zu{ }ſpät. Ich danke Ihnen beſtens für ihre freundliche
               Abſicht und grüße ihre Frau\pwindex{Schnitzler, Olga 17.\,1.\,1882 Wien – 13.\,1.\,1970 Lugano@\textsc{Schnitzler, Olga} (17.\,1.\,1882 Wien – 13.\,1.\,1970 Lugano), \emph{Schauspielerin, Sängerin}|pwv} und Sie herzlichſt.\pend
           
\pstart
           Ihr{\\[\baselineskip]}\spacefill\mbox{Gustav}\pend
           \leftskip=0em{}\selectlanguage{ngerman}\endnumbering\briefempfaengerindex{Schnitzler, Arthur@\textsc{Schnitzler, Arthur}!zzzSchwarzkopf, Gustav@\emph{von Gustav Schwarzkopf}!1908-03-062@{6. 3. 1908}|)be}\mylabel{L04270h}
\begin{anhang}
\end{anhang}\newcommand{\dateiname}{L04270}\newcommand{\titel}{Gustav Schwarzkopf an Arthur Schnitzler, 6. 3. 1908}\newcommand{\editorInnen}{Herausgegeben von Jahnke, SelmaMüller, Martin Anton}\input{../tex-inputs/latex-leseansicht-abspann}
